\documentclass[11pt]{article}
\usepackage[margin=0.9in]{geometry}
\usepackage{amsmath,amssymb}
\usepackage{xcolor}
\usepackage{textcomp}
\usepackage{fancyhdr}
\usepackage[hidelinks]{hyperref}
\definecolor{accent}{HTML}{1F6F8B}
\definecolor{soft}{HTML}{555555}
\definecolor{boxbg}{HTML}{F4F8FA}
\setlength{\parindent}{0pt}
\setlength{\parskip}{4pt}
\pagestyle{fancy}\fancyhf{}
\renewcommand{\headrulewidth}{0.4pt}
\lhead{\small\color{soft}FE Environmental \textemdash\ PEFEPrep}
\rhead{\small\color{soft}\rightmark}
\cfoot{\small\color{soft}\thepage}
\newcommand{\correct}[1]{\textbf{#1}\;{\color{accent}$\checkmark$}}
\newcommand{\optline}[2]{\par\noindent\hspace{1.2em}(#1)\hspace{0.6em}#2}

\begin{document}
\begin{center}
{\Huge\bfseries\color{accent} Fundamental Principles}\\[4pt]
{\large FE Environmental \textemdash\ Day 10}\\[2pt]
{\color{soft} Study set \textbullet\ 2026-06-26 \textbullet\ 18 questions \textbullet\ every numeric answer recomputed in Python}
\end{center}
\vspace{6pt}\hrule\vspace{10pt}
\markright{Fundamental Principles}
\par\noindent\colorbox{boxbg}{\parbox{\dimexpr\linewidth-2\fboxsep\relax}{%
\vspace{2pt}{\small\textbf{\color{accent}Handbook fidelity.}\ All formulas confirmed against NCEES FE Reference Handbook v10.6. Population projection equations (p. 334): Linear Pt = P0 + k∆t, Geometric Pt = P0e\textasciicircum{}(k∆t), Percent Pt = P0(1+k)\textasciicircum{}n, Ratio-Correlation P2/P2R = P1/P1R, Decreasing-Rate-of-Increase Pt = P0+(S$-$P0)(1$-$e\textasciicircum{}($-$k(t$-$t0))) --- all five models verbatim. Steady-State Reactor Parameters table (p. 330--331): CMFR and PFR first-order, zero-order, and second-order retention times and effluent concentrations verbatim. Engineering stress σ = F/A0, true stress σT = F/A, Hooke's Law σ = E$\varepsilon$ (p. 121). Thermal expansion coefficient $\varepsilon$ = $\alpha$∆T (p. 127 Materials Science). Thermal deformations δt = $\alpha$L(T$-$To) (p. 131 Mechanics of Materials). Fire hydrant rated capacity QR = QF$\times$(HR/HF)\textasciicircum{}0.54 and discharge Q = 29.8D\textsuperscript{2}CdP\textasciicircum{}0.5 (p. 304--305, Civil Engineering, NFPA 291). NOT in Handbook v10.6 (supplemental): the term 'logistic growth' (Handbook calls it Decreasing-Rate-of-Increase Growth); 'arithmetic growth' (Handbook calls it Linear Projection = Algebraic Projection); 'geometric' appears only in the synonym list 'Log Growth = Exponential Growth = Geometric Growth' --- all questions use Handbook nomenclature. The CMFR 2nd-order effluent formula Ct = ($\surd$(4kC0$\theta$+1)$-$1)/(2k$\theta$) is derived by solving the CMFR steady-state equation; the Handbook table gives the $\theta$ formula, from which Ct is algebraically equivalent --- tagged as handbook. True stress formula σT = σ$\times$(1+$\varepsilon$) for constant-volume plastic deformation is derivable from Handbook definitions (σT = F/A, σ = F/A0) but the combined form is not explicitly printed; Q14 gives both areas explicitly so students use σT = F/A directly (handbook). For USCS Q13: E = 29$\times$10\textsuperscript{6} psi is a standard steel value not tabulated in Handbook --- provided in the stem (tagged in-stem). Ratio-correlation: k = P2R/P1R (the constant) is rearranged from the Handbook form; the formula itself is in Handbook (tagged handbook).}\vspace{2pt}}}
\vspace{10pt}
\filbreak
\textbf{\color{accent}Q1.}\quad {\small\color{soft}[D10-Q01 \textbullet\ MCQ \textbullet\ Population projection --- linear (algebraic) model]}\par\vspace{2pt}
A small city had a population of 25,000 in 2010. Using the linear (algebraic) projection model, the population in 2020 is most nearly:

Linear Projection: Pt = P0 + k∆t

The growth rate k = 800 persons/year (determined from historical data).\par\vspace{3pt}
{\small\color{soft}Relevant:}\ $P_t = P_0 + k\,\Delta t$\par\vspace{3pt}
\optline{A}{30,000}
\optline{B}{31,000}
\optline{C}{\correct{33,000}}
\optline{D}{35,000}
\par\vspace{3pt}
\vspace{2pt}\par\noindent\colorbox{boxbg}{\parbox{\dimexpr\linewidth-2\fboxsep\relax}{%
\vspace{2pt}\textbf{Answer:} (C)\par\vspace{2pt}
{\small Using the linear projection formula: $P_t = P_0 + k\Delta t = 25{,}000 + 800\times10 = 25{,}000 + 8{,}000 = 33{,}000$.

Option (A) 30,000: uses k = 500 persons/yr (too low). Option (B) 31,000: uses k = 600 persons/yr. Option (D) 35,000: uses k = 1,000 persons/yr (too high). The linear model assumes a constant annual increment --- no compounding.}\par\vspace{2pt}
{\footnotesize\color{soft}Handbook: Environmental Engineering --- Population Projection Equations, Linear Projection = Algebraic Projection: Pt = P0 + k∆t (p. 334) \textbullet\ Handbook p. 334 --- Population Projection Equations}\vspace{2pt}
}}
\vspace{9pt}
\filbreak
\textbf{\color{accent}Q2.}\quad {\small\color{soft}[D10-Q02 \textbullet\ MCQ \textbullet\ Population projection --- geometric (log/exponential) model]}\par\vspace{2pt}
Census data for a municipality show a population of 50,000 in 2010 and 65,000 in 2020. Assuming geometric (log) growth:

Log Growth = Exponential Growth = Geometric Growth: Pt = P0e\textasciicircum{}(k∆t)

The projected population in 2025 is most nearly:\par\vspace{3pt}
{\small\color{soft}Relevant:}\ $P_t = P_0 e^{k\Delta t}$ \quad $k = \dfrac{\ln(P_t / P_0)}{\Delta t}$\par\vspace{3pt}
\optline{A}{70,000}
\optline{B}{72,000}
\optline{C}{\correct{74,000}}
\optline{D}{76,000}
\par\vspace{3pt}
\vspace{2pt}\par\noindent\colorbox{boxbg}{\parbox{\dimexpr\linewidth-2\fboxsep\relax}{%
\vspace{2pt}\textbf{Answer:} (C)\par\vspace{2pt}
{\small Step 1 --- find k from the two census points (2010$\rightarrow$2020, ∆t = 10 yr):
\[k = \frac{\ln(65{,}000/50{,}000)}{10} = \frac{\ln(1.30)}{10} = \frac{0.2624}{10} = 0.02624\text{ yr}^{-1}\]

Step 2 --- project from 2020 to 2025 (∆t = 5 yr from 2020 base):
\[P_{2025} = 65{,}000 \times e^{0.02624\times5} = 65{,}000 \times e^{0.1312} = 65{,}000 \times 1.1401 = 74{,}107\]

Option (A) 70,000: misapplies k, uses only additive estimate. Option (B) 72,000: underestimates k. Option (D) 76,000: overestimates. The geometric model compounds continuously --- larger populations grow faster each year.}\par\vspace{2pt}
{\footnotesize\color{soft}Handbook: Environmental Engineering --- Population Projection Equations, Log Growth = Exponential Growth = Geometric Growth: Pt = P0e\textasciicircum{}(k∆t) (p. 334) \textbullet\ Handbook p. 334 --- Population Projection Equations}\vspace{2pt}
}}
\vspace{9pt}
\filbreak
\textbf{\color{accent}Q3.}\quad {\small\color{soft}[D10-Q03 \textbullet\ MCQ \textbullet\ Population projection --- percent (discrete) growth model]}\par\vspace{2pt}
A growing town has a current population of 80,000. It is expected to grow at 2.5\% per year. Using the percent growth model, the population after 8 years is most nearly:

Percent Growth: Pt = P0(1 + k)\textasciicircum{}n\par\vspace{3pt}
{\small\color{soft}Relevant:}\ $P_t = P_0(1 + k)^n$\par\vspace{3pt}
\optline{A}{92,000}
\optline{B}{97,000}
\optline{C}{\correct{97,500}}
\optline{D}{102,000}
\par\vspace{3pt}
\vspace{2pt}\par\noindent\colorbox{boxbg}{\parbox{\dimexpr\linewidth-2\fboxsep\relax}{%
\vspace{2pt}\textbf{Answer:} (C)\par\vspace{2pt}
{\small $P_t = 80{,}000 \times (1 + 0.025)^8 = 80{,}000 \times (1.025)^8$.

$(1.025)^8$: $(1.025)^2 = 1.0506$; $(1.025)^4 = 1.1038$; $(1.025)^8 = 1.2184$.

$P_t = 80{,}000 \times 1.2184 = 97{,}472 \approx 97{,}500$.

Option (A) 92,000: uses n = 5 instead of 8. Option (B) 97,000: rounds k incorrectly. Option (D) 102,000: misapplies k = 3\%. The percent growth model is discrete (annual compounding), unlike the geometric model which compounds continuously.}\par\vspace{2pt}
{\footnotesize\color{soft}Handbook: Environmental Engineering --- Population Projection Equations, Percent Growth: Pt = P0(1+k)\textasciicircum{}n (p. 334) \textbullet\ Handbook p. 334 --- Population Projection Equations}\vspace{2pt}
}}
\vspace{9pt}
\filbreak
\textbf{\color{accent}Q4.}\quad {\small\color{soft}[D10-Q04 \textbullet\ MCQ \textbullet\ Population projection --- ratio-and-correlation method]}\par\vspace{2pt}
At the last census a city had a population of 40,000 and its surrounding metropolitan region had 500,000. In the new forecast, the region is projected to reach 620,000. Using the ratio-and-correlation method, the projected city population is most nearly:

Ratio and Correlation Growth: P2/P2R = P1/P1R\par\vspace{3pt}
{\small\color{soft}Relevant:}\ $\dfrac{P_2}{P_{2R}} = \dfrac{P_1}{P_{1R}}$\par\vspace{3pt}
\optline{A}{44,800}
\optline{B}{46,500}
\optline{C}{\correct{49,600}}
\optline{D}{52,000}
\par\vspace{3pt}
\vspace{2pt}\par\noindent\colorbox{boxbg}{\parbox{\dimexpr\linewidth-2\fboxsep\relax}{%
\vspace{2pt}\textbf{Answer:} (C)\par\vspace{2pt}
{\small The ratio-correlation method assumes the city's share of the regional population stays constant:
\[P_2 = P_{2R} \times \frac{P_1}{P_{1R}} = 620{,}000 \times \frac{40{,}000}{500{,}000} = 620{,}000 \times 0.080 = 49{,}600\]

Option (A) 44,800: divides by 620 instead of multiplying by the ratio. Option (B) 46,500: uses ratio of 40/500 but applies to wrong base. Option (D) 52,000: inflates the ratio. The method scales the city in proportion to its region --- useful when regional forecasts are more reliable than local ones.}\par\vspace{2pt}
{\footnotesize\color{soft}Handbook: Environmental Engineering --- Population Projection Equations, Ratio and Correlation Growth: P2/P2R = P1/P1R (p. 334) \textbullet\ Handbook p. 334 --- Population Projection Equations}\vspace{2pt}
}}
\vspace{9pt}
\filbreak
\textbf{\color{accent}Q5.}\quad {\small\color{soft}[D10-Q05 \textbullet\ MCQ \textbullet\ Population projection --- decreasing-rate-of-increase (saturation) model]}\par\vspace{2pt}
A city currently has a population of 30,000. Analysis of historical data indicates a saturation population S = 80,000 and a growth rate constant k = 0.10 yr\textsuperscript{--}\textsuperscript{1}. Using the decreasing-rate-of-increase model, the population 15 years from now is most nearly:

Decreasing-Rate-of-Increase Growth: Pt = P0 + (S $-$ P0)(1 $-$ e\textasciicircum{}($-$k(t$-$t0)))\par\vspace{3pt}
{\small\color{soft}Relevant:}\ $P_t = P_0 + (S - P_0)\left(1 - e^{-k(t-t_0)}\right)$\par\vspace{3pt}
\optline{A}{58,000}
\optline{B}{64,000}
\optline{C}{\correct{68,800}}
\optline{D}{76,000}
\par\vspace{3pt}
\vspace{2pt}\par\noindent\colorbox{boxbg}{\parbox{\dimexpr\linewidth-2\fboxsep\relax}{%
\vspace{2pt}\textbf{Answer:} (C)\par\vspace{2pt}
{\small $P_t = 30{,}000 + (80{,}000 - 30{,}000)(1 - e^{-0.10 \times 15})$
$= 30{,}000 + 50{,}000\,(1 - e^{-1.5})$
$= 30{,}000 + 50{,}000\,(1 - 0.2231)$
$= 30{,}000 + 50{,}000 \times 0.7769$
$= 30{,}000 + 38{,}843 = 68{,}843 \approx 68{,}800$.

Option (A) 58,000: uses t = 10 yr instead of 15. Option (B) 64,000: uses k = 0.07. Option (D) 76,000: assumes nearly complete saturation (overestimates). This model captures the idea that growth slows as the population approaches its carrying capacity S.}\par\vspace{2pt}
{\footnotesize\color{soft}Handbook: Environmental Engineering --- Population Projection Equations, Decreasing-Rate-of-Increase Growth: Pt = P0 + (S$-$P0)(1$-$e\textasciicircum{}($-$k(t$-$t0))) (p. 334) \textbullet\ Handbook p. 334 --- Population Projection Equations}\vspace{2pt}
}}
\vspace{9pt}
\filbreak
\textbf{\color{accent}Q6.}\quad {\small\color{soft}[D10-Q06 \textbullet\ MCQ \textbullet\ CMFR --- first-order reaction, retention time]}\par\vspace{2pt}
A completely mixed flow reactor (CMFR) is used to treat an influent with concentration C0 = 200 mg/L. The treatment target is an effluent concentration Ct = 20 mg/L. The first-order rate constant is k = 0.5 d\textsuperscript{--}\textsuperscript{1}. The required mean retention time $\theta$ is most nearly:

For a CMFR with first-order decay (r = $-$kC): $\theta$ = (C0/Ct $-$ 1)/k\par\vspace{3pt}
{\small\color{soft}Relevant:}\ $\theta = \dfrac{C_0/C_t - 1}{k}$\par\vspace{3pt}
\optline{A}{4.6 d}
\optline{B}{9.0 d}
\optline{C}{\correct{18.0 d}}
\optline{D}{36.0 d}
\par\vspace{3pt}
\vspace{2pt}\par\noindent\colorbox{boxbg}{\parbox{\dimexpr\linewidth-2\fboxsep\relax}{%
\vspace{2pt}\textbf{Answer:} (C)\par\vspace{2pt}
{\small $\theta = \dfrac{C_0/C_t - 1}{k} = \dfrac{200/20 - 1}{0.5} = \dfrac{10 - 1}{0.5} = \dfrac{9}{0.5} = 18.0$ d.

Option (A) 4.6 d: this is the PFR retention time for the same removal (see Q07) --- wrong reactor model applied. Option (B) 9.0 d: forgets to divide by k. Option (D) 36.0 d: doubles the answer (arithmetic error on the ratio). A CMFR is always less efficient than a PFR for the same kinetics because it operates at the effluent (lowest) concentration throughout.}\par\vspace{2pt}
{\footnotesize\color{soft}Handbook: Environmental Engineering --- Steady-State Reactor Parameters table, Ideal CMFR, First-Order: $\theta$ = (C0/Ct $-$ 1)/k (p. 330--331) \textbullet\ Handbook p. 330--331 --- Steady-State Reactor Parameters (Constant Density Systems)}\vspace{2pt}
}}
\vspace{9pt}
\filbreak
\textbf{\color{accent}Q7.}\quad {\small\color{soft}[D10-Q07 \textbullet\ MCQ \textbullet\ Plug-flow reactor (PFR) --- first-order reaction, retention time]}\par\vspace{2pt}
A plug-flow reactor (PFR) must achieve 90\% removal of a contaminant with C0 = 200 mg/L and Ct = 20 mg/L. The first-order rate constant k = 0.5 d\textsuperscript{--}\textsuperscript{1}. The required retention time $\theta$ is most nearly:

For an ideal PFR with first-order decay: $\theta$ = ln(C0/Ct)/k\par\vspace{3pt}
{\small\color{soft}Relevant:}\ $\theta = \dfrac{\ln(C_0/C_t)}{k}$\par\vspace{3pt}
\optline{A}{2.3 d}
\optline{B}{\correct{4.6 d}}
\optline{C}{9.0 d}
\optline{D}{18.0 d}
\par\vspace{3pt}
\vspace{2pt}\par\noindent\colorbox{boxbg}{\parbox{\dimexpr\linewidth-2\fboxsep\relax}{%
\vspace{2pt}\textbf{Answer:} (B)\par\vspace{2pt}
{\small $\theta = \dfrac{\ln(200/20)}{0.5} = \dfrac{\ln(10)}{0.5} = \dfrac{2.303}{0.5} = 4.61$ d $\approx 4.6$ d.

Option (A) 2.3 d: ln(10)/1.0 --- forgets to divide by k or uses k = 1 d\textsuperscript{--}\textsuperscript{1}. Option (C) 9.0 d: applies CMFR numerator formula (C0/Ct$-$1 = 9) to PFR --- wrong model. Option (D) 18.0 d: CMFR result. Comparison: the PFR achieves 90\% removal in 4.6 d vs the CMFR's 18.0 d --- a 3.9$\times$ advantage because the PFR maintains high driving-force concentration along the plug length.}\par\vspace{2pt}
{\footnotesize\color{soft}Handbook: Environmental Engineering --- Steady-State Reactor Parameters table, Ideal Plug Flow, First-Order: $\theta$ = ln(C0/Ct)/k (p. 330--331) \textbullet\ Handbook p. 330--331 --- Steady-State Reactor Parameters (Constant Density Systems)}\vspace{2pt}
}}
\vspace{9pt}
\filbreak
\textbf{\color{accent}Q8.}\quad {\small\color{soft}[D10-Q08 \textbullet\ MCQ \textbullet\ CMFR --- zero-order reaction, retention time]}\par\vspace{2pt}
A CMFR is operated with a zero-order reaction rate k = 10 mg/(L$\cdot$d). The influent concentration is C0 = 100 mg/L and the target effluent is Ct = 40 mg/L. The required mean retention time $\theta$ is most nearly:

For a CMFR (and ideal batch and PFR) with zero-order decay (r = $-$k): $\theta$ = (C0 $-$ Ct)/k\par\vspace{3pt}
{\small\color{soft}Relevant:}\ $\theta = \dfrac{C_0 - C_t}{k}$\par\vspace{3pt}
\optline{A}{4.0 d}
\optline{B}{\correct{6.0 d}}
\optline{C}{10.0 d}
\optline{D}{14.0 d}
\par\vspace{3pt}
\vspace{2pt}\par\noindent\colorbox{boxbg}{\parbox{\dimexpr\linewidth-2\fboxsep\relax}{%
\vspace{2pt}\textbf{Answer:} (B)\par\vspace{2pt}
{\small $\theta = \dfrac{C_0 - C_t}{k} = \dfrac{100 - 40}{10} = \dfrac{60}{10} = 6.0$ d.

Option (A) 4.0 d: uses C0 $-$ Ct = 40 instead of 60 (confuses Ct and the reduction). Option (C) 10.0 d: uses C0/k = 100/10 (ignores that we only need to reduce to 40, not to zero). Option (D) 14.0 d: arithmetic error. Note: for zero-order kinetics the CMFR, PFR, and batch all have the same retention time --- unlike first-order kinetics where the PFR is more efficient.}\par\vspace{2pt}
{\footnotesize\color{soft}Handbook: Environmental Engineering --- Steady-State Reactor Parameters table, Ideal CMFR, Zero-Order: $\theta$ = (C0$-$Ct)/k (p. 330--331) \textbullet\ Handbook p. 330--331 --- Steady-State Reactor Parameters (Constant Density Systems)}\vspace{2pt}
}}
\vspace{9pt}
\filbreak
\textbf{\color{accent}Q9.}\quad {\small\color{soft}[D10-Q09 \textbullet\ Numeric \textbullet\ CMFR --- second-order reaction, effluent concentration]}\par\vspace{2pt}
A CMFR operates with a second-order decay reaction (r = $-$kC\textsuperscript{2}) with rate constant k = 0.02 L/(mg$\cdot$d) and a mean retention time $\theta$ = 5 d. The influent concentration is C0 = 50 mg/L. The effluent concentration Ct is most nearly:

For a CMFR with second-order decay: $C_t = \dfrac{\sqrt{4k C_0 \theta + 1} - 1}{2k\theta}$\par\vspace{3pt}
{\small\color{soft}Relevant:}\ $C_t = \dfrac{\sqrt{4k C_0 \theta + 1} - 1}{2k\theta}$\par\vspace{3pt}
{\small\color{soft}Numeric entry.}\par\vspace{3pt}
\vspace{2pt}\par\noindent\colorbox{boxbg}{\parbox{\dimexpr\linewidth-2\fboxsep\relax}{%
\vspace{2pt}\textbf{Answer:} 1\par\vspace{2pt}
{\small $C_t = \dfrac{\sqrt{4(0.02)(50)(5)+1}-1}{2(0.02)(5)} = \dfrac{\sqrt{21}-1}{0.2} = \dfrac{4.583-1}{0.2} = \dfrac{3.583}{0.2} = 17.9$ mg/L.

Step by step: $4kC_0\theta = 4(0.02)(50)(5) = 20$; inner radical: $\sqrt{21} = 4.583$; numerator: $4.583-1=3.583$; denominator: $2(0.02)(5)=0.2$; $C_t = 3.583/0.2 = 17.9$ mg/L.

Option (A) 14.6 mg/L: uses incorrect radical expansion. Option (C) 25.0 mg/L: uses first-order formula $C_0/(1+k\theta)$ with wrong k units. Option (D) 35.7 mg/L: neglects the subtraction of 1 in the numerator. The second-order CMFR requires solving a quadratic steady-state balance, yielding the formula above.}\par\vspace{2pt}
{\footnotesize\color{soft}Handbook: Environmental Engineering --- Steady-State Reactor Parameters table, Ideal CMFR, Second-Order: Ct formula (p. 330--331) \textbullet\ Handbook p. 330--331 --- Steady-State Reactor Parameters (Constant Density Systems)}\vspace{2pt}
}}
\vspace{9pt}
\filbreak
\textbf{\color{accent}Q10.}\quad {\small\color{soft}[D10-Q10 \textbullet\ MCQ \textbullet\ CMFR vs PFR --- comparison for first-order removal]}\par\vspace{2pt}
A treatment process must achieve 90\% removal of a contaminant that degrades via first-order kinetics. Two reactor configurations are considered: an ideal CMFR and an ideal PFR, each with the same first-order rate constant k. Which statement best compares the required hydraulic retention times ($\theta$)?\par\vspace{3pt}
{\small\color{soft}Relevant:}\ $\theta_{\text{CMFR}} = \dfrac{C_0/C_t - 1}{k}$ \quad $\theta_{\text{PFR}} = \dfrac{\ln(C_0/C_t)}{k}$\par\vspace{3pt}
\optline{A}{$\theta$\_CMFR < $\theta$\_PFR (CMFR is more efficient)}
\optline{B}{$\theta$\_CMFR $\approx$ $\theta$\_PFR (both reactors perform equally)}
\optline{C}{\correct{$\theta$\_CMFR $\approx$ 3.9 $\times$ $\theta$\_PFR (CMFR requires \textasciitilde{}4$\times$ longer HRT)}}
\optline{D}{$\theta$\_CMFR = 9 $\times$ $\theta$\_PFR (CMFR requires exactly 9$\times$ longer HRT)}
\par\vspace{3pt}
\vspace{2pt}\par\noindent\colorbox{boxbg}{\parbox{\dimexpr\linewidth-2\fboxsep\relax}{%
\vspace{2pt}\textbf{Answer:} (C)\par\vspace{2pt}
{\small For 90\% removal: $C_0/C_t = 10$.

$\theta_{\text{CMFR}} = \dfrac{10-1}{k} = \dfrac{9}{k}$

$\theta_{\text{PFR}} = \dfrac{\ln 10}{k} = \dfrac{2.303}{k}$

Ratio: $\theta_{\text{CMFR}}/\theta_{\text{PFR}} = 9/2.303 = 3.91$

Option (A): wrong --- the CMFR is less efficient, not more, because it operates at the (lowest) effluent concentration throughout the entire volume. Option (B): wrong --- they differ significantly for first-order kinetics. Option (D): $\theta_{\text{CMFR}} = 9/k$ and $\theta_{\text{PFR}} = \ln(10)/k \approx 2.3/k$, so the ratio is 3.9, not 9. The PFR advantage exists because it maintains a high concentration (high reaction rate) along the length of the reactor.}\par\vspace{2pt}
{\footnotesize\color{soft}Handbook: Environmental Engineering --- Steady-State Reactor Parameters table, Ideal CMFR and Ideal Plug Flow, First-Order (p. 330--331) \textbullet\ Handbook p. 330--331 --- Steady-State Reactor Parameters (Constant Density Systems)}\vspace{2pt}
}}
\vspace{9pt}
\filbreak
\textbf{\color{accent}Q11.}\quad {\small\color{soft}[D10-Q11 \textbullet\ MCQ \textbullet\ Engineering stress --- axial loading (SI)]}\par\vspace{2pt}
A steel rod with an initial diameter of 20 mm carries an axial tensile load of 40 kN. The engineering stress in the rod is most nearly:

Engineering stress: σ = F/A0\par\vspace{3pt}
{\small\color{soft}Relevant:}\ $\sigma = \dfrac{F}{A_0}$ \quad $A_0 = \dfrac{\pi d^2}{4}$\par\vspace{3pt}
\optline{A}{63.7 MPa}
\optline{B}{100 MPa}
\optline{C}{\correct{127 MPa}}
\optline{D}{159 MPa}
\par\vspace{3pt}
\vspace{2pt}\par\noindent\colorbox{boxbg}{\parbox{\dimexpr\linewidth-2\fboxsep\relax}{%
\vspace{2pt}\textbf{Answer:} (C)\par\vspace{2pt}
{\small Initial cross-sectional area:
\[A_0 = \frac{\pi (0.020)^2}{4} = \frac{\pi \times 4 \times 10^{-4}}{4} = 3.142 \times 10^{-4}\text{ m}^2\]

Engineering stress:
\[\sigma = \frac{F}{A_0} = \frac{40{,}000\text{ N}}{3.142 \times 10^{-4}\text{ m}^2} = 127.3 \times 10^6\text{ Pa} \approx 127\text{ MPa}\]

Option (A) 63.7 MPa: uses diameter instead of radius in computing area (A = $\pi$D\textsuperscript{2}/4 not $\pi$D\textsuperscript{2}). Option (B) 100 MPa: error in unit conversion. Option (D) 159 MPa: uses r = 20 mm (full diameter as radius). Engineering stress always uses the original (undeformed) cross-sectional area A0.}\par\vspace{2pt}
{\footnotesize\color{soft}Handbook: Materials Science/Structure of Matter --- Engineering stress: σ = F/A0, where A0 = initial cross-sectional area (p. 121) \textbullet\ Handbook p. 121 --- Materials Science/Structure of Matter: Stress and Strain}\vspace{2pt}
}}
\vspace{9pt}
\filbreak
\textbf{\color{accent}Q12.}\quad {\small\color{soft}[D10-Q12 \textbullet\ MCQ \textbullet\ Young's modulus --- Hooke's Law (SI)]}\par\vspace{2pt}
A specimen of aluminum alloy is subjected to a uniaxial tensile stress of 140 MPa, producing an elastic strain of 0.002 (dimensionless). The elastic modulus (Young's modulus) of the material is most nearly:

Hooke's Law: σ = E$\varepsilon$\par\vspace{3pt}
{\small\color{soft}Relevant:}\ $E = \dfrac{\sigma}{\varepsilon}$\par\vspace{3pt}
\optline{A}{28 GPa}
\optline{B}{56 GPa}
\optline{C}{\correct{70 GPa}}
\optline{D}{140 GPa}
\par\vspace{3pt}
\vspace{2pt}\par\noindent\colorbox{boxbg}{\parbox{\dimexpr\linewidth-2\fboxsep\relax}{%
\vspace{2pt}\textbf{Answer:} (C)\par\vspace{2pt}
{\small \[E = \frac{\sigma}{\varepsilon} = \frac{140 \times 10^6\text{ Pa}}{0.002} = 70 \times 10^9\text{ Pa} = 70\text{ GPa}\]

This is consistent with the elastic modulus of aluminum (\textasciitilde{}70 GPa).

Option (A) 28 GPa: divides by 5 instead of 0.002 (misplaces decimal). Option (B) 56 GPa: arithmetic error. Option (D) 140 GPa: uses E = σ/$\varepsilon$ but forgets to apply $\varepsilon$ (divides by 1 instead of 0.002). Hooke's Law applies only in the elastic (linear) region of the stress-strain curve.}\par\vspace{2pt}
{\footnotesize\color{soft}Handbook: Materials Science/Structure of Matter --- Hooke's Law: σ = E$\varepsilon$, where E = elastic modulus (also called modulus, modulus of elasticity, Young's modulus) (p. 121) \textbullet\ Handbook p. 121 --- Materials Science/Structure of Matter: Elastic modulus (Young's modulus)}\vspace{2pt}
}}
\vspace{9pt}
\filbreak
\textbf{\color{accent}Q13.}\quad {\small\color{soft}[D10-Q13 \textbullet\ MCQ \textbullet\ Hooke's Law --- elastic strain from applied stress (USCS)]}\par\vspace{2pt}
A structural steel member carries an axial stress of 29,000 psi. The steel has an elastic modulus E = 29 $\times$ 10\textsuperscript{6} psi (given). The elastic strain is most nearly:

Hooke's Law: σ = E$\varepsilon$\par\vspace{3pt}
{\small\color{soft}Relevant:}\ $\varepsilon = \dfrac{\sigma}{E}$\par\vspace{3pt}
\optline{A}{0.0005}
\optline{B}{\correct{0.0010}}
\optline{C}{0.0015}
\optline{D}{0.0020}
\par\vspace{3pt}
\vspace{2pt}\par\noindent\colorbox{boxbg}{\parbox{\dimexpr\linewidth-2\fboxsep\relax}{%
\vspace{2pt}\textbf{Answer:} (B)\par\vspace{2pt}
{\small \[\varepsilon = \frac{\sigma}{E} = \frac{29{,}000\text{ psi}}{29 \times 10^6\text{ psi}} = 1.0 \times 10^{-3} = 0.0010\]

Option (A) 0.0005: divides by 58 $\times$ 10\textsuperscript{6} (doubles E). Option (C) 0.0015: arithmetic error. Option (D) 0.0020: uses E = 14.5 $\times$ 10\textsuperscript{6} psi (half of correct value). At $\varepsilon$ = 0.001, this steel is comfortably in the elastic range (yield strain for structural steel $\approx$ 0.001--0.002).}\par\vspace{2pt}
{\footnotesize\color{soft}Handbook: Materials Science/Structure of Matter --- Hooke's Law: σ = E$\varepsilon$ (p. 121). E = 29$\times$10\textsuperscript{6} psi for steel is provided in the stem. \textbullet\ Handbook p. 121 --- Hooke's Law σ = E$\varepsilon$; E value provided in stem}\vspace{2pt}
}}
\vspace{9pt}
\filbreak
\textbf{\color{accent}Q14.}\quad {\small\color{soft}[D10-Q14 \textbullet\ MCQ \textbullet\ True stress vs engineering stress]}\par\vspace{2pt}
A tensile test specimen has an initial cross-sectional area A0 = 200 mm\textsuperscript{2} and is loaded with F = 30,000 N. At a point during plastic deformation, the actual (current) cross-sectional area has necked to A = 170 mm\textsuperscript{2}. The true stress at this point is most nearly:

True stress: σT = F/A  |  Engineering stress: σ = F/A0\par\vspace{3pt}
{\small\color{soft}Relevant:}\ $\sigma_T = \dfrac{F}{A}$ \quad $\sigma = \dfrac{F}{A_0}$\par\vspace{3pt}
\optline{A}{150.0 MPa (engineering stress --- uses initial area)}
\optline{B}{164.7 MPa}
\optline{C}{\correct{176.5 MPa (true stress --- uses actual area)}}
\optline{D}{188.2 MPa}
\par\vspace{3pt}
\vspace{2pt}\par\noindent\colorbox{boxbg}{\parbox{\dimexpr\linewidth-2\fboxsep\relax}{%
\vspace{2pt}\textbf{Answer:} (C)\par\vspace{2pt}
{\small True stress uses the actual (current) cross-sectional area:
\[\sigma_T = \frac{F}{A} = \frac{30{,}000\text{ N}}{170 \times 10^{-6}\text{ m}^2} = 176.5 \times 10^6\text{ Pa} = 176.5\text{ MPa}\]

Engineering stress (for comparison):
\[\sigma = \frac{F}{A_0} = \frac{30{,}000}{200 \times 10^{-6}} = 150.0\text{ MPa}\]

Option (A) 150 MPa: this is the engineering stress --- uses undeformed A0. Option (B) 164.7 MPa: uses an intermediate area estimate. Option (D) 188.2 MPa: uses A = 160 mm\textsuperscript{2} (too small). True stress exceeds engineering stress after yielding because the area decreases. Before yielding they are essentially equal.}\par\vspace{2pt}
{\footnotesize\color{soft}Handbook: Materials Science/Structure of Matter --- True stress: σT = F/A (actual area); Engineering stress: σ = F/A0 (initial area). 'True stress is load divided by actual cross-sectional area whereas engineering stress is load divided by the initial area.' (p. 121) \textbullet\ Handbook p. 121 --- Materials Science/Structure of Matter: True stress and Engineering stress}\vspace{2pt}
}}
\vspace{9pt}
\filbreak
\textbf{\color{accent}Q15.}\quad {\small\color{soft}[D10-Q15 \textbullet\ MCQ \textbullet\ Thermal deformation --- axial member (SI)]}\par\vspace{2pt}
A structural steel member is 4.0 m long. It is heated from an initial temperature T0 = 20\textdegree{}C to a final temperature T = 80\textdegree{}C. The thermal expansion coefficient of steel is $\alpha$ = 12 $\times$ 10\textsuperscript{--}\textsuperscript{6} /\textdegree{}C. The thermal deformation δt is most nearly:

Thermal Deformations: δt = $\alpha$L(T $-$ T0)\par\vspace{3pt}
{\small\color{soft}Relevant:}\ $\delta_t = \alpha L (T - T_0)$\par\vspace{3pt}
\optline{A}{1.4 mm}
\optline{B}{\correct{2.9 mm}}
\optline{C}{3.6 mm}
\optline{D}{5.0 mm}
\par\vspace{3pt}
\vspace{2pt}\par\noindent\colorbox{boxbg}{\parbox{\dimexpr\linewidth-2\fboxsep\relax}{%
\vspace{2pt}\textbf{Answer:} (B)\par\vspace{2pt}
{\small \[\delta_t = \alpha L(T - T_0) = (12 \times 10^{-6}\text{ /^{\circ}C})(4.0\text{ m})(80 - 20)\text{ ^{\circ}C}\]
\[= 12 \times 10^{-6} \times 4.0 \times 60 = 2{,}880 \times 10^{-6}\text{ m} = 2.88\text{ mm} \approx 2.9\text{ mm}\]

Option (A) 1.4 mm: uses ∆T = 30\textdegree{}C (misreads ∆T = T $-$ T0 as T/2). Option (C) 3.6 mm: uses ∆T = 75\textdegree{}C or L = 5 m. Option (D) 5.0 mm: uses $\alpha$ = 21 $\times$ 10\textsuperscript{--}\textsuperscript{6} (incorrect for steel). The formula is in the Mechanics of Materials section; δt is positive for expansion (T > T0).}\par\vspace{2pt}
{\footnotesize\color{soft}Handbook: Mechanics of Materials --- Thermal Deformations: δt = $\alpha$L(T $-$ To), where δt = deformation caused by a change in temperature, $\alpha$ = temperature coefficient of expansion (p. 131) \textbullet\ Handbook p. 131 --- Mechanics of Materials: Thermal Deformations}\vspace{2pt}
}}
\vspace{9pt}
\filbreak
\textbf{\color{accent}Q16.}\quad {\small\color{soft}[D10-Q16 \textbullet\ MCQ \textbullet\ Thermal expansion coefficient --- strain from temperature change]}\par\vspace{2pt}
The thermal expansion coefficient of an aluminum alloy is $\alpha$ = 23 $\times$ 10\textsuperscript{--}\textsuperscript{6} /\textdegree{}C. A component is heated by ∆T = 50\textdegree{}C. The resulting thermal engineering strain $\varepsilon$ is most nearly:

Thermal Properties: $\varepsilon$ = $\alpha$∆T\par\vspace{3pt}
{\small\color{soft}Relevant:}\ $\varepsilon = \alpha\,\Delta T$\par\vspace{3pt}
\optline{A}{0.00092}
\optline{B}{\correct{0.00115}}
\optline{C}{0.00138}
\optline{D}{0.00160}
\par\vspace{3pt}
\vspace{2pt}\par\noindent\colorbox{boxbg}{\parbox{\dimexpr\linewidth-2\fboxsep\relax}{%
\vspace{2pt}\textbf{Answer:} (B)\par\vspace{2pt}
{\small \[\varepsilon = \alpha\,\Delta T = (23 \times 10^{-6}\text{ /^{\circ}C})(50\text{ ^{\circ}C}) = 1{,}150 \times 10^{-6} = 1.15 \times 10^{-3} = 0.00115\]

Option (A) 0.00092: uses $\alpha$ = 18.4 $\times$ 10\textsuperscript{--}\textsuperscript{6} (value for a different material). Option (C) 0.00138: uses ∆T = 60\textdegree{}C. Option (D) 0.00160: uses $\alpha$ = 32 $\times$ 10\textsuperscript{--}\textsuperscript{6}. This formula gives the free thermal strain --- the strain that would occur in an unconstrained member. If the member is restrained, a thermal stress σ = E$\varepsilon$ = E$\alpha$∆T develops.}\par\vspace{2pt}
{\footnotesize\color{soft}Handbook: Materials Science/Structure of Matter --- Thermal Properties: 'The thermal expansion coefficient is the ratio of engineering strain to the change in temperature.' $\varepsilon$ = $\alpha$∆T (p. 127) \textbullet\ Handbook p. 127 --- Materials Science: Thermal Properties (thermal expansion coefficient)}\vspace{2pt}
}}
\vspace{9pt}
\filbreak
\textbf{\color{accent}Q17.}\quad {\small\color{soft}[D10-Q17 \textbullet\ MCQ \textbullet\ Fire hydrant --- rated capacity at 20 psi residual]}\par\vspace{2pt}
During a fire hydrant flow test, the static pressure is PS = 80 psi and the residual pressure is PR = 45 psi at a total test flow of QF = 1,200 gpm. The rated capacity QR at 20 psi residual is most nearly:

Fire Hydrant Rated Capacity (NFPA 291):
QR = QF $\times$ (HR/HF)\textasciicircum{}0.54
where HR = PS $-$ 20 and HF = PS $-$ PR\par\vspace{3pt}
{\small\color{soft}Relevant:}\ $Q_R = Q_F \times \left(\dfrac{H_R}{H_F}\right)^{0.54}$\par\vspace{3pt}
\optline{A}{1,200 gpm}
\optline{B}{1,400 gpm}
\optline{C}{\correct{1,600 gpm}}
\optline{D}{1,800 gpm}
\par\vspace{3pt}
\vspace{2pt}\par\noindent\colorbox{boxbg}{\parbox{\dimexpr\linewidth-2\fboxsep\relax}{%
\vspace{2pt}\textbf{Answer:} (C)\par\vspace{2pt}
{\small \[H_R = P_S - 20 = 80 - 20 = 60\text{ psi}\]
\[H_F = P_S - P_R = 80 - 45 = 35\text{ psi}\]
\[Q_R = 1{,}200 \times \left(\frac{60}{35}\right)^{0.54} = 1{,}200 \times (1.714)^{0.54}\]

$(1.714)^{0.54}$: $\ln(1.714) = 0.538$; $0.54 \times 0.538 = 0.291$; $e^{0.291} = 1.338$.

\[Q_R = 1{,}200 \times 1.338 = 1{,}605\text{ gpm} \approx 1{,}600\text{ gpm}\]

Option (A) 1,200 gpm: the test flow --- no extrapolation applied. Option (B) 1,400 gpm: underestimates the exponent or drops the 0.54. Option (D) 1,800 gpm: uses exponent 1.0 (linear extrapolation) instead of 0.54.}\par\vspace{2pt}
{\footnotesize\color{soft}Handbook: Civil Engineering --- Formula for Calculating Rated Capacity at 20 psi from Fire Hydrant: QR = QF $\times$ (HR/HF)\textasciicircum{}0.54, Source: NFPA Standard 291 (p. 304) \textbullet\ Handbook p. 304 --- Civil Engineering: Fire Hydrant Rated Capacity (NFPA 291)}\vspace{2pt}
}}
\vspace{9pt}
\filbreak
\textbf{\color{accent}Q18.}\quad {\small\color{soft}[D10-Q18 \textbullet\ MCQ \textbullet\ Fire hydrant --- discharge to atmosphere]}\par\vspace{2pt}
A fire hydrant with a smooth, well-rounded outlet (Cd = 0.90) has an outlet diameter of 2.5 in. A pitot gauge reading at the nozzle indicates P = 36 psi. The flow rate discharged to the atmosphere is most nearly:

Fire Hydrant Discharging to Atmosphere (NFPA 291):
Q = 29.8 D\textsuperscript{2} Cd P\textasciicircum{}(1/2)\par\vspace{3pt}
{\small\color{soft}Relevant:}\ $Q = 29.8\,D^2\,C_d\,P^{1/2}$\par\vspace{3pt}
\optline{A}{670 gpm}
\optline{B}{840 gpm}
\optline{C}{\correct{1,000 gpm}}
\optline{D}{1,340 gpm}
\par\vspace{3pt}
\vspace{2pt}\par\noindent\colorbox{boxbg}{\parbox{\dimexpr\linewidth-2\fboxsep\relax}{%
\vspace{2pt}\textbf{Answer:} (C)\par\vspace{2pt}
{\small \[Q = 29.8 \times (2.5)^2 \times 0.90 \times \sqrt{36}\]
\[= 29.8 \times 6.25 \times 0.90 \times 6.0\]
\[= 29.8 \times 6.25 = 186.25\]
\[\times\, 0.90 = 167.6\]
\[\times\, 6.0 = 1{,}006\text{ gpm} \approx 1{,}000\text{ gpm}\]

Option (A) 670 gpm: uses Cd = 0.60 (wrong coefficient) or misses the D\textsuperscript{2} squaring. Option (B) 840 gpm: uses Cd = 0.75 (intermediate coefficient for a square-edged outlet). Option (D) 1,340 gpm: omits Cd (treats Cd = 1.0). The smooth, well-rounded outlet has Cd = 0.90; square-edged = 0.80; projecting into barrel = 0.70 (per NFPA 291 table in Handbook).}\par\vspace{2pt}
{\footnotesize\color{soft}Handbook: Civil Engineering --- Fire Hydrant Discharging to Atmosphere: Q = 29.8D\textsuperscript{2}CdP\textasciicircum{}(1/2), Cd = 0.90 for smooth, well-rounded outlet (p. 304--305, NFPA 291) \textbullet\ Handbook p. 304--305 --- Civil Engineering: Fire Hydrant Discharging to Atmosphere (NFPA 291)}\vspace{2pt}
}}
\vspace{9pt}
\end{document}